\chapter{Decisões de Design e Acessibilidade}
\label{chap:design}

\section{Paleta de Cores}
\label{chap4:sec:cores}

A interface usa um tema escuro com um acento dourado/âmbar. Não foi uma escolha por moda --- cada cor tem uma razão de ser.

O fundo usa \texttt{\#0e0c0a}, um preto com subtom castanho-quente em vez do preto puro. Em ecrãs OLED, o contraste extremo entre fundo preto e texto branco pode ser desconfortável em sessões longas; o subtom quente atenua isso. Visualmente também remete para o ambiente de um clube de jazz, com iluminação ténue, o que se enquadra bem no tema da aplicação.

O acento dourado (\texttt{\#c9a84c}) é a cor de todos os elementos interativos. A associação é direta: saxofone, trompete, trombone --- os instrumentos centrais do jazz têm exatamente esta cor de latão. Além disso o dourado sobre fundo escuro cumpre os rácios de contraste do \ac{WCAG}~2.1 nível AA~\cite{wcag21}, o que não é coincidência. O texto principal usa \texttt{\#f0ebe0}, um branco ligeiramente creme, e o secundário \texttt{\#a09080}. Ambos mantêm a coerência quente da paleta e criam hierarquia clara sem precisar de tamanhos de letra muito diferentes.

Para estados de erro e sucesso escolhemos deliberadamente um vermelho suave (\texttt{\#d95555}) e um verde-azulado (\texttt{\#5cb88a}) em vez de verde puro. A razão é simples: verde e vermelho são as cores mais problemáticas para pessoas com daltonismo, e o verde-azulado já cria distância percetiva suficiente para reduzir confusões mesmo sem ativar nenhum modo de acessibilidade.

\section{Tipografia}
\label{chap4:sec:tipografia}

Usámos exclusivamente a família \textbf{Inter} em pesos de 300 a 700. O protótipo da Parte~1 combinava Playfair Display nos títulos com Inter no corpo, mas abandonámos essa combinação por três razões práticas.

Primeiro, as serifas da Playfair ficavam com aliasing visível a densidades típicas de smartphones Android. Segundo, dois sistemas tipográficos num espaço compacto criam atrito visual sem benefício real --- a hierarquia que a Playfair pretendia introduzir consegue-se perfeitamente com Inter Bold a um tamanho ligeiramente maior. Terceiro, fontes sans-serif com formas abertas são mais fáceis de ler para utilizadores com dislexia~\cite{bda2023}, e a Inter foi desenhada especificamente para interfaces digitais, com aberturas largas nos caracteres e altura-x elevada.

\begin{table}[!htb]
\centering\small
\begin{tabular}{llll}
\toprule
\textbf{Elemento} & \textbf{Peso} & \textbf{Tamanho} & \textbf{Uso} \\
\midrule
Títulos de ecrã   & 700 Bold     & 20--22 px & Cabeçalho de cada ecrã \\
Títulos de secção & 600 SemiBold & 17--18 px & Títulos de cards e painéis \\
Corpo             & 400 Regular  & 14--15 px & Texto corrido \\
Labels            & 500--600     & 11--13 px & Categorias, metadados \\
Texto secundário  & 300--400     & 11--13 px & Datas, dicas \\
\bottomrule
\end{tabular}
\caption{Hierarquia tipográfica Inter na interface do GenJazz.}
\label{tab:tipografia}
\end{table}

\section{Acessibilidade para Utilizadores com Daltonismo}
\label{chap4:sec:daltonismo}

A aplicação inclui três modos de acessibilidade cromática, acessíveis nas definições do perfil. As paletas foram baseadas na paleta de Wong~(2011)~\cite{wong2011}, desenvolvida para garantir discriminabilidade nos principais tipos de deficiência de visão cromática. Cada modo substitui apenas as variáveis CSS de acento, erro e sucesso --- os fundos e textos mantêm-se inalterados.

\begin{table}[!htb]
\centering\small
\begin{tabular}{lp{3.2cm}p{3.2cm}p{3.2cm}}
\toprule
\textbf{Papel} & \textbf{Deuteranopia} & \textbf{Protanopia} & \textbf{Tritanopia} \\
\midrule
Acento  & \texttt{\#56B4E9} azul céu    & \texttt{\#56B4E9} azul céu    & \texttt{\#D55E00} vermelhão \\
Erro    & \texttt{\#E69F00} laranja      & \texttt{\#E69F00} laranja      & \texttt{\#CC3311} verm.-laranja \\
Sucesso & \texttt{\#0072B2} azul escuro  & \texttt{\#009E73} verde-azulado & \texttt{\#009E73} verde-azulado \\
\bottomrule
\end{tabular}
\caption{Paletas de acessibilidade para os três tipos de daltonismo suportados.}
\label{tab:daltonismo}
\end{table}

A deuteranopia --- dificuldade em perceber o verde, presente em cerca de 5--6\% da população masculina~\cite{birch2012} --- foi o caso que mais justificou esta funcionalidade, uma vez que o dourado da paleta padrão tem uma componente verde que o torna difícil de distinguir neste modo de visão. Para a protanopia (ausência de sensibilidade ao vermelho, \textasciitilde1\%) a paleta é parecida mas com o sucesso ajustado, pois protanopes conseguem distinguir verdes deslocados para o ciano. Para a tritanopia (afeta o azul, \textasciitilde0,01\%) o dourado é substituído por vermelhão, que não depende da distinção azul/amarelo.

\begin{figure}[!htb]
\centering
\begin{subfigure}[t]{0.28\textwidth}
  \includegraphics[width=\textwidth]{04_home.png}
  \caption{Padrão.}
\end{subfigure}\hfill
\begin{subfigure}[t]{0.28\textwidth}
  \includegraphics[width=\textwidth]{11_home_deuteranopia.png}
  \caption{Deuteranopia.}
\end{subfigure}\hfill
\begin{subfigure}[t]{0.28\textwidth}
  \includegraphics[width=\textwidth]{12_generation_tritanopia.png}
  \caption{Tritanopia.}
\end{subfigure}
\caption{Comparação entre o modo padrão e dois modos de acessibilidade cromática.}
\label{fig:daltonismo}
\end{figure}

A implementação é feita exclusivamente via variáveis CSS: ao mudar o modo, adiciona-se o atributo \texttt{data-colorblind} ao elemento \texttt{<html>} e as variáveis corretas são automaticamente aplicadas em toda a interface, sem alterar nenhum componente individualmente.

\section{Outras Considerações de Acessibilidade}
\label{chap4:sec:outras}

Para além dos modos de daltonismo, tomámos algumas decisões que melhoram a acessibilidade geral. Todos os pares de cor texto/fundo foram verificados para cumprir o nível AA do \ac{WCAG}~2.1. Os botões e elementos tácteis têm pelo menos 44$\times$44~px, seguindo as recomendações da Apple e do W3C. A hierarquia da interface é transmitida por tamanho, peso e posição, não dependendo só de cor --- o que faz com que a interface funcione bem mesmo no modo de alto contraste do sistema operativo.
